\documentclass[10pt,a4paper]{amsart}
\usepackage[margin=19mm]{geometry}
\usepackage{amsmath,amssymb,mathtools}
\usepackage[scr=rsfs,cal=euler]{mathalfa}
\usepackage{xfrac}
\usepackage[unicode,hidelinks,bookmarks=false]{hyperref}
\newcommand{\ie}{i.e.\ }
\newcommand{\cf}{cf.\ }
\newcommand{\resp}{respectively\ }
\newcommand{\Subsection}{\subsection}
\providecommand{\nobreakdash}{\nobreak\hbox{-}}
\setlength{\emergencystretch}{2em}
\multlinegap0pt
\usepackage{xcolor}
\usepackage{amssymb}
\usepackage{cancel}
\newtheorem{proposition}{Proposition}
\newtheorem{theorem}{Theorem}
\title{Uniqueness of capillary disks in three-dimensional domains}
\author{Henrique Nogueira Bastos}
\date{Source: arXiv:2603.25929v1 (CC BY 4.0)}
\begin{document}
\maketitle

\noindent\emph{Source and licence.} Uniqueness of capillary disks in three-dimensional domains. Version arXiv:2603.25929v1 (CC BY 4.0), distributed by arXiv under the Creative Commons Attribution 4.0 International licence, http://creativecommons.org/licenses/by/4.0/, as recorded in the arXiv licence metadata for this exact version and shown on the abstract page https://arxiv.org/abs/2603.25929v1. Full licence notice: \texttt{licenses/arxiv-2603.25929-CC-BY-4.0.txt}.\par\smallskip

\noindent\textit{Editorial scope.} Counted unit: the article's theorem mainA (source0.tex lines 252--254). The article's complete original proof of it is delivered with it (source0.tex lines 255--357, one \texttt{proof} environment of 103 lines). No mathematics is written, repaired or re-derived: the statement and the proof are the article's own lines, in the article's own notation, and no retained line is substituted or re-flowed.  Retained uncounted as the source-local context the proof needs: lines 148--156; lines 152--154; lines 240--244.  Nothing was omitted from any retained span, so the package carries no omission record.  No retained line contains a citation, so the wrapper carries no bibliography and no bibliography entry.  No retained line contains a figure environment, an image-inclusion command or drawing code, so no image is delivered and the package declares no asset dependency.  Editorial formatting only, in addition: an AMS \texttt{amsart} preamble that restates the article's own declarations and macros used by the retained lines, namely the environments \texttt{proposition}, \texttt{theorem} on separate counters, together with the standard packages those lines need.  The retained environments print in this document as Proposition 1, Theorem 1 in order of appearance, whereas the article prints the same items with the numbering of its own full document.  The complete native source of the article as distributed in the arXiv e-print for this exact version is bundled unchanged as \texttt{sources/197253/source0.tex}.
\begin{proposition}\label{prop:condição de bordo}
    Let $\Sigma$ be an immersed capillary surface in $\Omega$ with contact angle $\alpha\in(0,\pi)$. Let $(u,v,z)$ be adapted coordinates  to $(p,T_p\Sigma)$ for some $p\in \partial(\Sigma\cap  \Omega)$, chosen so that $\{u=0\}$ corresponds to $\partial \Omega$ and $(0,v,z)$ are isothermal coordinates on $\partial \Omega$, with conformal factor $\lambda^2>0$. 
    
    If $\Sigma$ is locally parametrized as $X(u,v)=(u,v,h(u,v))$ near $p$, then
    \begin{equation}\label{eq:prop2.1}
     -g_{uu}(0,v,h(0,v))(1+h_v^2(0,v))\cos^2\alpha+\lambda^2(0,v,h(0,v)) h_u(0,v)^2\sin^2\alpha=0
    \end{equation}
    for all $v$ sufficiently close to $0$.
\end{proposition}

    \begin{equation}
     -g_{uu}(0,v,h(0,v))(1+h_v^2(0,v))\cos^2\alpha+\lambda^2(0,v,h(0,v)) h_u(0,v)^2\sin^2\alpha=0
    \end{equation}

\begin{equation}\label{formula sigma}
\sigma=g(\partial_z(0),N(0)) \begin{bmatrix}
    w_{uu} & w_{uv} \\ w_{vu} & w_{vv}
\end{bmatrix}+o((\sqrt{u^2+v^2})^{k-2}),
\end{equation}

\begin{theorem}\label{thm: mainA}
     Outside its singular set, the principal lines of $\sigma$ on $\Sigma$ are either tangential to $\partial(\Sigma\cap \Omega)$ or intersect it orthogonally.
\end{theorem}

\begin{proof} To study the behavior of the line fields at a point $p\in \partial (\Sigma\cap\Omega)$, we choose adapted coordinates $(u,v,z)$ around $p$ as the one chosen in Proposition \ref{prop:condição de bordo}.

%such that $p=(0,0,0)$ and the set $\{u=0\}$ corresponding to $\partial\Omega$ around $p$ with $\partial_u$ being a normal vector field along $\partial\Omega$ and $g(\partial_u(p),\partial_v(p))=0$ and $g(\partial_z(p),\partial_v(p))=0$. Let $c(t)=(u(t),v(t),z(t))$ be a parametrization of $\partial(\Sigma\cap \Omega)$ in a neighborhood of $p$.

We consider the parametrizations $X(u,v)=(u,v,h(u,v))$ for $\Sigma$ and $\widetilde{X}(u,v)=(u,v,\widetilde{h}(u,v))$ for the comparison surface $S(p,T_p\Sigma)=:\widetilde{\Sigma}$ in the transitive family, with $p=(0,0,0)$ and $h_v(0,0)=\widetilde{h}_v(0,0)=0$.

We first analyze the free boundary case, that is, when the contact angle $\alpha=\pi/2$. Substituting $\alpha=\pi/2$ into \eqref{eq:prop2.1}, we obtain $h_u(0,v)=0$, and hence $h_{uv}(0,0)=0$. %is orthogonal to $\partial \Omega$ along the boundary of $\Sigma\cap \Omega$, so that
%\begin{equation*}
 %   0 = g(X_u,\partial_z)(c(t)) = g_{uz}(c(t))+h_u(0,v(t))g_{zz}(c(t))
%\end{equation*}
%along $c(t)$. Since $c(0)=p$ and $c'(0)=\partial_v(p)$, and $h_u(0,0)=0$, we obtain
%\begin{equation*}
    %0 = g_{uz,v}(p)+h_{uv}(0,0)g_{zz}(p).
%\end{equation*}
Taking the difference with the analogous identity for $\widetilde{\Sigma}$, we obtain
\begin{equation*}
    w_{uv}(0,0)=(h-\widetilde{h})_{uv}(0,0)=0
\end{equation*}
where $w$ is the homogeneous harmonic polynomial appearing in \eqref{formula sigma}. Therefore, the directions $X_u(0,0)=\partial_u(p)$ and $X_v(0,0)=\partial_v(p)$ are principal directions of $\sigma$. In particular, $\partial_v(p)$ is not an asymptotic direction of $\sigma$.  The Theorem is proven in this case.

We now turn to the non free boundary case, that is, constant contact angle $\alpha\neq \pi/2$. Let $\nu_{\Sigma\cap \Omega}$ denote the conormal vector of $\Sigma\cap\Omega$, and $N_{\partial\Omega}$ denote the normal vector to $\partial\Omega$. Keeping the same notation from the proof of Proposition \ref{prop:condição de bordo}, we write
    \begin{align*}
        \nu_{\Sigma\cap \Omega}& =\frac{1}{\sqrt{G^{uu}}}(G^{uu}X_u+G^{uv}X_v) \\
        N_{\partial \Omega}& =\frac{1}{\sqrt{B^{uu}}}(B^{uu}X_u+B^{uv}X_v+B^{uz}\partial_z),
    \end{align*}
where $B_{ij}=g(X_i,X_j)$. (Here we abuse notation by writing $X_z=\partial_z$).

The capillary condition implies that the vectors $\nu(t)=\nu_{\Sigma\cap\Omega}(c(t))$ and $N(t)=N_{\partial \Omega}(c(t))$ satisfy
\begin{align*}
    \cos(\alpha) = & g(\nu(t),N(t)) \\
    = & \frac{1}{\sqrt{G^{uu}}}[G^{uu}g(X_u,N)+G^{uv}g(X_v,N)] \\
    = & \sqrt{G^{uu}}g(X_u,N)= \frac{\sqrt{G^{uu}}}{\sqrt{B^{uu}}}.
\end{align*}
Differentiating with respect to $t$, we obtain
\begin{align*}
    0 & = \frac{\sqrt{B^{uu}}}{2\sqrt{G^{uu}}}
    \ \frac{(G^{uu})_{,v}B^{uu}-(B^{uu})_{,v}G^{uu}}{(B^{uu})^2} \\
    & = \frac{\sqrt{B^{uu}}}{2\sqrt{G^{uu}}}\ \frac{B^{ui}B^{uj}B_{ij,v}G^{uu}-G^{ui}G^{uj}G_{ij,v}B^{uu}}{(B^{uu})^2},
\end{align*}
and hence
\begin{equation}\label{eq:derivada ao longo de c}
B^{ui}B^{uj}B_{ij,v}G^{uu}-G^{ui}G^{uj}G_{ij,v}B^{uu} = 0.
\end{equation}
Recall that $\partial_u$ is normal to $\partial\Omega$ and that $(0,v,z)$ are isothermal coordinates on $\partial\Omega$ in our choice of adapted coordinates. Also, $h_v(0,0)=\widetilde{h}_v(0,0)=0$. Since
\begin{equation*}    
    G(p)=\begin{bmatrix}
            G_{uu} & 0 \\
            0 & G_{vv}
        \end{bmatrix} \text{ and }
    B(p)=\begin{bmatrix}
            G_{uu} & 0 & B_{uz} \\
            0 & G_{vv} & 0 \\
            B_{uz} & 0 & B_{zz}
        \end{bmatrix},
\end{equation*}
and the inverse of this matrix at $p$ 
\begin{equation*}
    B(p)^{-1}=\frac{1}{G_{uu}B_{zz}-B_{uz}^2}\begin{bmatrix}
            B_{zz} & 0 & -B_{uz} \\
            0 & (G_{uu}B_{zz}-B_{uz}^2)/G_{vv} & 0 \\
            -B_{zu} & 0 & B_{uu}
        \end{bmatrix},
\end{equation*}
then, at $p$, equation \eqref{eq:derivada ao longo de c} yields
\begin{align*}
    0 = & (B^{uu})^2B_{uu,v}+2B^{uu}B^{uz}B_{uz,v}+(B^{uz})^2B_{zz,v}-G^{uu}B^{uu}G_{uu,v} \\
    = & B^{uu}G_{uu,v}(B^{uu}-G^{uu})+B^{uz}(2B^{uu}B_{uz,v}+B^{uz}B_{zz,v}) \\
    = & B^{uu}G_{uu,v}\left(\frac{B_{zz}}{G_{uu}B_{zz}-B_{uz}^2}-\frac{1}{G_{uu}}\right)+B^{uz}(2B^{uu}B_{uz,v}+B^{uz}B_{zz,v}) \\
    = & B^{uu}G_{uu,v}\frac{B_{uz}^2}{(G_{uu}B_{zz}-B_{uz}^2)G_{uu}}+B^{uz}(2B^{uu}B_{uz,v}+B^{uz}B_{zz,v}) \\
    = & -B^{uu}G_{uu,v}\frac{B_{uz}B^{uz}}{G_{uu}}+B^{uz}(2B^{uu}B_{uz,v}+B^{uz}B_{zz,v}).
\end{align*}
Taking the difference between the above expressions evaluated on $\Sigma$ and $\widetilde{\Sigma}$ at $p$, we obtain
\begin{align}\label{eq:auxiliar}
    0%= %& [-B^{uu}G_{uu,v}\frac{B_{uz}B^{uz}}{G_{uu}}+B^{uz}(2B^{uu}B_{uz,v}+B^{uz}B_{zz,v})] \nonumber \\ 
    %& -[-B^{uu}\widetilde{G}_{uu,v}\frac{B_{uz}B^{uz}}{G_{uu}}+B^{uz}(2B^{uu}\widetilde{B}_{uz,v}+B^{uz}\widetilde{B}_{zz,v})] \nonumber \\ 
    = & -B^{uu}\frac{B_{uz}B^{uz}}{G_{uu}}(G_{uu}-\widetilde{G}_{uu})_{,v} + 2B^{uu}B^{uz}(B_{uz}-\widetilde{B}_{uz})_{,v},
\end{align}
where we used that $B_{zz}=g(\partial_z,\partial_z)\circ c(t)$, $\widetilde{B}_{zz}=g(\partial_z,\partial_z)\circ \widetilde{c}(t)$, so that $B_{zz,v}=\widetilde{B}_{zz,v}$ at $t=0$, since $c(0)=\widetilde{c}(0)=p$ and $c'(0)=\widetilde{c}'(0)=\partial_v(p)$. 

Recall that $h_u(0,0)-\widetilde{h}_u(0,0)=0$, since the tangent planes of $\Sigma$ and $\widetilde{\Sigma}$ coincide at $p=(0,0,0)$. Since
\begin{equation*}
    G_{uu}=g_{uu}+2h_u g_{uz} + h^2_u g_{zz} \text{ and } \widetilde{G}_{uu}=g_{uu}+2\widetilde{h}_u g_{uz} + \widetilde{h}^2_u g_{zz},
\end{equation*}
and similarly
\begin{equation*}
    B_{uz}=g_{uz}+h_u g_{zz} \text{ and } \widetilde{B}_{uz}=g_{uz}+\widetilde{h}_u g_{zz},
\end{equation*}
differentiating with respect to $v$ at $p$ and substituting into \eqref{eq:auxiliar}, we obtain
\begin{align*}
    0= & -2B_{uz}B^{uz}(h-\widetilde{h})_{uv}B_{uz}+2B^{uz}G_{uu}(h-\widetilde{h})_{uv} B_{zz} \\
    = & 2B^{uz}(h-\widetilde{h})_{uv}(B_{zz}G_{uu}-B_{uz}^2).
\end{align*}
Note that $B_{zz}G_{uu}-B^2_{uz}\neq 0$. Therefore, at $(0,0)$,
\begin{equation*}
    (h-\widetilde{h})_{uv}=0 \text{ or } B^{uz}=0. 
\end{equation*}

Since we are now analyzing the case of constant angle $\alpha\neq \pi/2$, the second possibility cannot occur. Thus
\begin{equation*}
    (h-\widetilde{h})_{uv}(0,0)=0 
\end{equation*}
and the conclusion follows as in the free boundary case.    
\end{proof}

\end{document}
